\begin{exercisechunk}
\item \textbf{Representing parity with one ReLU layer.}
\label[exercise]{ex:l7-parity-representation}
\exerciseprofile{2}{2}{2}{\exproof\quad\excalculation}{%
Verify that a width-$(|S|+1)$ ReLU network represents a fixed parity exactly.}
Fix $S\subseteq[d]$, write $k=|S|$, and let $\mathbf 1_S\in\{0,1\}^d$ be
the indicator vector of $S$. Consider
\begin{equation*}
N(x)=[-\mathbf 1_S^\top x+k+1]_+
+\sum_{j=1}^k4j(-1)^j
[-\mathbf 1_S^\top x+k+1-2j]_+.
\end{equation*}
Prove that
\[
N(x)=\chi_S(x)
\qquad\text{for every }x\in\{-1,1\}^d.
\]
Conclude that the network in \cref{prop:parity-representation} needs at most
$k+1$ hidden units.

\begin{hints}
Let $a$ be the number of negative coordinates of $x$ indexed by $S$. First
show that $-\mathbf 1_S^\top x+k+1=2a+1$ and identify which ReLU terms are
active. If $\widetilde N(a)$ denotes the resulting value, compare
$\widetilde N(a)$ with $\widetilde N(a-1)$. Pair consecutive terms when
evaluating $\sum_{j=1}^{a-1}j(-1)^j$.
\end{hints}
\end{exercisechunk}
